% =====================================================================
\section{Visi\'on y alcance}\label{sec:vision}
% =====================================================================

\subsection{Declaraci\'on de visi\'on}
Para la persona responsable de JC\textbar{}cell, que necesita encontrar con facilidad los datos
de los equipos recibidos y conocer el avance de cada reparaci\'on, se propone una herramienta
sencilla que concentre el registro de ingreso, el diagn\'ostico y el estado del servicio. Se busca
reducir el tiempo que hoy dedica a revisar anotaciones y ayudarle a identificar los trabajos
pendientes, manteniendo su atenci\'on directa, sus diagn\'osticos y sus reparaciones. Esta es una
propuesta acad\'emica basada en SE-01 (B2, B4, C3--C4, D3; EV-04), no un compromiso de construir o
instalar el sistema; su alcance debe validarse con el responsable del taller.

\subsection{Objetivos de negocio}
El participante expres\'o que le gustar\'ia encontrar antes la informaci\'on de cada equipo, ubicar
las reparaciones pendientes y evitar confusiones y cobros olvidados (SE-01, C4, D1; EV-04).
La entrevista no incluy\'o mediciones iniciales ni metas acordadas. Por eso, los indicadores de la
Tabla~\ref{tab:objetivos} son propuestas para medir esos resultados; las l\'ineas base, metas y
plazos se completar\'an tras revisar ejemplos anonimizados y acordarlos con el responsable. Hasta
entonces, los objetivos siguen como borrador y no se consideran verificables.

\begin{table}[H]
  \centering
  \caption{Objetivos de negocio e indicadores propuestos para validar.}
  \label{tab:objetivos}
  \small
  \renewcommand{\arraystretch}{1.25}
  \noindent\begin{tabularx}{\linewidth}{@{}L{1.1cm}YL{2.6cm}L{2.4cm}L{2.4cm}@{}}
    \toprule
    \textbf{ID} & \textbf{Resultado buscado} & \textbf{Indicador} & \textbf{L\'inea base} & \textbf{Meta y plazo} \\
    \midrule
    ON-01 & Reducir el tiempo que toma localizar los datos y el estado de un servicio. Fuente: C3--C4, D1. & Mediana de minutos de las búsquedas completadas; registrar también las fallidas. & \pendiente{medir con ejemplos anonimizados (E3)} & \pendiente{acordar meta y fecha tras medir} \\
    ON-02 & Mejorar la consistencia del seguimiento de las reparaciones. Fuente: B4, C2, C4. & Servicios con estado coincidente / servicios verificables de la muestra $\times 100$. & \pendiente{revisar ejemplos de trabajos (E3)} & \pendiente{acordar porcentaje y plazo} \\
    ON-03 & Disminuir los cobros pendientes que quedan sin seguimiento al entregar un equipo. Fuente: B7, C2, D1. & Número de entregas con saldo pendiente confirmado, sin marca de pendiente en el registro. & \pendiente{revisar ejemplos de cobro (E3)} & \pendiente{acordar meta y fecha tras medir} \\
    \bottomrule
  \end{tabularx}
  \normalsize
\end{table}

Las fuentes de cada objetivo corresponden a preguntas de SE-01, recogidas en EV-04 y en el
Anexo~\ref{anx:registro-entrevista}. Los indicadores son propuestas del equipo, no mediciones
solicitadas o aprobadas por el participante.

Para medirlos, se propone acordar con el responsable el período, el tamaño de la muestra y los
criterios de selección antes de recoger los datos:
\begin{itemize}[nosep]
  \item \textbf{ON-01:} medir desde que se solicita localizar un servicio hasta encontrar sus
        datos y estado. Acordar un tiempo máximo y registrar las búsquedas fallidas por separado;
        la mediana se interpretará junto con la cantidad de búsquedas completadas y fallidas.
  \item \textbf{ON-02:} comparar el estado anotado con el que pueda reconstruirse junto al
        responsable, usando registros de trabajo o entrega del mismo momento de referencia.
        Los casos sin respaldo suficiente se contarán aparte, sin asumir que coinciden.
  \begin{samepage}
  \item \textbf{ON-03:} contrastar las notas del servicio y la entrega con la aclaración del
        responsable para identificar saldos pendientes que no quedaron marcados. Registrar solo
        el estado del cobro, sin datos de clientes; si no puede verificarse, dejar el caso pendiente.
  \par\end{samepage}
\end{itemize}
Las comparaciones posteriores usarán los mismos criterios; en ON-03 se usarán muestras del mismo
tamaño y se indicará el período revisado. La medición y sus condiciones aún deben acordarse y
realizarse; por eso se mantienen pendientes las líneas base, metas y fechas.

\subsection{Alcance}
La propuesta se centra en que el responsable pueda registrar y consultar los servicios del taller.
Las capacidades de la Tabla~\ref{tab:alcance} se derivan de SE-01 (B2--B7, C4, D1--D3, E2;
EV-04) y siguen pendientes de revisi\'on con el responsable; la entrevista no defini\'o todos los
detalles de operaci\'on.

\begin{table}[H]
  \centering
  \caption{Alcance propuesto para validar con el responsable.}
  \label{tab:alcance}
  \small
  \renewcommand{\arraystretch}{1.2}
  \noindent\begin{tabularx}{\linewidth}{@{}L{1.3cm}YL{1.6cm}L{2cm}@{}}
    \toprule
    \textbf{ID} & \textbf{Capacidad propuesta} & \textbf{Objetivo} & \textbf{Fuente} \\
    \midrule
    \multicolumn{4}{@{}l}{\textit{Dentro del alcance propuesto}} \\
    AL-01 & Registrar la recepci\'on del equipo, sus datos b\'asicos, la falla reportada y el contacto asociado. & ON-02 & B2, C4 \\
    AL-02 & Guardar el diagn\'ostico, la autorizaci\'on del cliente y el estado de cada reparaci\'on. & ON-02 & B3--B4 \\
    AL-03 & Consultar y localizar equipos y servicios activos, incluidos los que esperan revisi\'on, autorizaci\'on, repuesto o entrega. & ON-01 & B4--B5, C3--C4 \\
    AL-04 & Anotar el valor comunicado y si el servicio qued\'o pagado o pendiente, sin procesar pagos. & ON-03 & B3, B7 \\
    AL-05 & Consultar el historial de reparaciones asociadas a un equipo. & ON-01 & E2 \\
    \midrule
    \multicolumn{4}{@{}l}{\textit{Exclusiones propuestas, pendientes de acuerdo}} \\
    FA-01 & Envío automático de mensajes o integración con WhatsApp; se describieron llamadas y mensajes, pero no se definió su automatización. & --- & B6 \\
    FA-02 & Control de inventario, reposición automática o pedidos a proveedores; se mencionaron faltantes y pérdidas, pero no se definieron las funciones de ese módulo. & --- & B5, C2 \\
    FA-03 & Procesamiento de pagos, facturación electrónica o contabilidad; se describió el registro del cobro, sin definir esas funciones. & --- & B7 \\
    FA-04 & Administración de reglas y plazos de garantía; se mencionó su seguimiento, pero no se especificaron las condiciones. & --- & B8, E2 \\
    \bottomrule
  \end{tabularx}
  \normalsize
\end{table}

La columna de fuente remite a preguntas de SE-01 (EV-04;
Anexo~\ref{anx:registro-entrevista}). Los vínculos con los objetivos son propuestas del equipo;
por ejemplo, B7 respalda el problema de cobros de ON-03 y la capacidad AL-04. AL-05 se relaciona
con ON-01 porque permitiría consultar antecedentes de un equipo; ese vínculo aún debe validarse.

Las exclusiones son propuestas del equipo, pendientes de acuerdo con el taller. La falta de
detalle en la entrevista no demuestra que una función sea innecesaria. Antes de cerrar el alcance
se confirmarán tanto las inclusiones como las exclusiones con el responsable.

\subsection{Restricciones conocidas}
La Tabla~\ref{tab:restricciones} separa las condiciones expresadas por el responsable de los datos
que todav\'ia deben confirmarse. No se fija una plataforma t\'ecnica ni una obligaci\'on legal que no
haya sido verificada para este proyecto.

\begin{table}[H]
  \centering
  \caption{Restricciones y asuntos por confirmar.}
  \label{tab:restricciones}
  \small
  \renewcommand{\arraystretch}{1.2}
  \noindent\begin{tabularx}{\linewidth}{@{}L{1.1cm}L{2.2cm}YL{2.4cm}@{}}
    \toprule
    \textbf{ID} & \textbf{Tipo} & \textbf{Condici\'on} & \textbf{Origen} \\
    \midrule
    RE-01 & Datos y privacidad & Los ejemplos usados en el proyecto deben mantenerse anonimizados; los originales y datos identificables no se publican. El tratamiento, acceso y conservaci\'on de datos reales del sistema a\'un deben definirse. & Manifiesto de evidencias; EV-03 y EV-04 \\
    RE-02 & Recursos & La propuesta debe considerar la computadora, los celulares y la conexi\'on disponibles, y no requerir equipos costosos ni conocimientos t\'ecnicos avanzados. La plataforma y la calidad m\'inima de conexi\'on quedan por validar. & SE-01, D2--D3; EV-04 \\
    RE-03 & Plazo del proyecto & La fecha l\'imite y el calendario de hitos deben confirmarse con el cronograma oficial antes de fijar metas. & Cronograma del PFC; fecha pendiente de confirmaci\'on \\
    \bottomrule
  \end{tabularx}
  \normalsize
\end{table}
